% Part: incompleteness
% Chapter: representability-in-q
% Section: c

\documentclass[../../../include/open-logic-section]{subfiles}

\begin{document}

\olfileid{inc}{req}{cee}
\olsection{$C$નાં વિધેયો}

$C$ એ નીચેના વિધેયો ધરાવતો સૌથી નાનો વિધેય-ગણ હોય:
\begin{enumerate}
\item $0$,
\item ઉત્તરગામી વિધેય,
\item સરવાળો,
\item ગુણાકાર,
\item પ્રક્ષેપો, અને
\item સમાનતા માટેનું લાક્ષણિક વિધેય, $\Char{=}$;
\end{enumerate}
અને જે નીચેની ક્રિયાઓ હેઠળ બંધ હોય:
\begin{enumerate}
\item સંયોજન, અને
\item નિયમિત વિધેયો પર લાગુ પડતી અનિયત મર્યાદા સુધીની શોધ.
\end{enumerate}
યાદ રાખો કે છેલ્લી મર્યાદાનો અર્થ એટલો જ છે કે
પરિણામી વિધેય સંપૂર્ણ હોય ત્યારે જ $\mu$ ક્રિયા
વાપરી શકાય. તેની સરખામણી \emph{સામાન્ય પુનરાવર્તી}
વિધેયોની વ્યાખ્યા સાથે કરો: અહીં આપણે સરવાળો, ગુણાકાર
અને $\Char{=}$ ઉમેર્યાં છે, પરંતુ આદિમ પુનરાવર્તન
દૂર કર્યું છે.

$C$માંનું દરેક વિધેય પુનરાવર્તી છે એ સ્પષ્ટ છે,
કારણ કે સરવાળો, ગુણાકાર અને $\Char{=}$ પણ
પુનરાવર્તી છે. વિપરીત દિશા પણ સાચી છે એ આપણે
બતાવીશું; એટલે $C$ની અન્ય સામગ્રી વડે આદિમ
પુનરાવર્તન કરી શકાય છે.

\end{document}
